\section{Introduction}

A \textbf{business process} (BP) is a coordinated sequence of activities an organisation performs to reach an outcome, and what makes it one is the coordination, not its size or complexity. \textbf{Business Process Management} (BPM) is the systematic design, analysis, and improvement of such processes: it automates workflow, orchestrates work across organisational units, reduces the risk of errors, provides real-time metrics, enforces deadlines, validates data, and cuts training costs. Software reached these goals late. Information systems started from data, so information modelling came first; the \textbf{workflow management systems} of the mid-nineties gave processes their first systematic support, but stayed confined to a few domains. BPM broadens that scope to \textbf{process-aware information systems} spanning multiple organisations, and aims at the \emph{correct realisation of business processes in software}.

That aim is pursued from three angles: formal methods prove properties, software development builds the systems, business administration chases cost and customer satisfaction. What holds them together is \textbf{abstraction}: each works on a shared \emph{model} rather than on the messy reality. Every organisation already runs processes; modelling makes them explicit enough to analyse, communicate, and change.

\begin{notebox}
  \emph{All models are wrong, but some are useful}, George Box.

  What a model discards may be irrelevant to the question at hand, and choosing what to leave out is the skill of modelling.
\end{notebox}

\subsection{Two directions of modelling}

A process model can be reached from either end. \textbf{BPMN} draws it top-down: an analyst designs the process by hand. Process mining goes the other way, deriving it bottom-up from recorded behaviour.

\textbf{BPMN} (Business Process Model and Notation) is a graphical language for describing processes: a collaboration, say a student and a supervisor during a thesis project, gives each participant its own swim-lane of activities, and the lanes exchange messages.

\textbf{Process mining} bridges process models and event data. A software system records events (messages, transactions, state changes) into \textbf{event logs}; three operations then relate logs to models:

\begin{itemize}
  \item \textbf{Discovery}: derive a process model from an event log with no prior model.
  \item \textbf{Conformance}: check how well an existing model matches the observed log.
  \item \textbf{Enhancement}: improve or extend an existing model using log data.
\end{itemize}

\subsection{Digression: Formal Modelling and the Importance of Proofs}

The Königsberg bridge problem shows why formal models and proofs matter in BPM. Königsberg's seven bridges joined two islands and two riverbanks, and a popular puzzle asked whether one could walk a circuit crossing each bridge exactly once.

\begin{wrapfigure}{r}{0.34\linewidth}
  \centering
  \vspace{-0.4cm}
  \begin{tikzpicture}[
    every node/.style={circle, draw=accent, thick, fill=accent!70,
      minimum size=9pt, inner sep=0pt},
    edge/.style={draw=accent, thick},
    lbl/.style={draw=none, fill=none}
  ]
    \node (C) at (0,0) {};
    \node (D) at (2.8,0) {};
    \node (A) at (1.4,0.95) {};
    \node (B) at (1.4,-0.95) {};
    % C-A double
    \draw[edge] (C) to[bend left=22] (A);
    \draw[edge] (C) to[bend right=22] (A);
    % C-B double
    \draw[edge] (C) to[bend left=22] (B);
    \draw[edge] (C) to[bend right=22] (B);
    % C-D single (central)
    \draw[edge] (C) to (D);
    % A-D single
    \draw[edge] (A) to (D);
    % B-D single
    \draw[edge] (B) to (D);
    \node[lbl, above=3pt] at (A.north) {\small\sffamily A};
    \node[lbl, below=3pt] at (B.south) {\small\sffamily B};
    \node[lbl, left=3pt] at (C.west) {\small\sffamily C};
    \node[lbl, right=3pt] at (D.east) {\small\sffamily D};
  \end{tikzpicture}
  \caption{The Königsberg graph.}
  \label{fig:euler-graph}
  \vspace{-0.4cm}
\end{wrapfigure}

Euler (1736) set about \emph{proving} no such walk exists. He abstracted the map to a graph: land areas became \textbf{nodes}, bridges became \textbf{edges}. A \textbf{path} is then a finite sequence of contiguous edges, and it is a \textbf{circuit} when it ends where it starts; either one is \textbf{Eulerian} if it uses every edge of the graph exactly once. The puzzle becomes a graph problem: given a connected graph, find an Eulerian circuit if one exists. The answer turns on the \textbf{degree} of each vertex, the number of edges attached to it, which is either \textbf{odd} or \textbf{even}.

\begin{theorem}[Euler's theorem]
  A connected graph $G$ contains an Eulerian circuit \textbf{if and only if} every vertex is even.
\end{theorem}

In the Königsberg graph all four vertices are odd (A, B, D have degree~3; C has degree~5), so the theorem immediately rules out an Eulerian circuit: no amount of trial and error can find a solution, because none exists.

\begin{theorem}[Euler's theorem, Eulerian path]
  A connected graph $G$ contains an Eulerian path \textbf{if and only if} $G$ has exactly $0$ or $2$ vertices of odd degree.

  When all vertices are even, the path is also a circuit. When exactly $2$ vertices are odd, they are the start and end of the path.
\end{theorem}

\textit{Proof (necessity).} Suppose $G$ contains an Eulerian path $P$. At every vertex $v$, the path enters and leaves the same number of times, except when $v$ is the start or end of $P$. If start and end are distinct, exactly two vertices are odd. If they coincide, no vertex is odd.\hfill$\square$

\textit{Proof (sufficiency).} Suppose $G$ has $0$ or $2$ odd vertices. If there are two odd vertices $u$ and $v$, add a temporary edge $(u, v)$: now every vertex is even, so by Euler's theorem the modified graph contains an Eulerian circuit $C$, and $C$ traverses $(u, v)$ at some point. Remove $(u, v)$ from $C$: what remains is an Eulerian path with $u$ and $v$ as endpoints. If there are no odd vertices, Euler's theorem applies directly and the circuit is itself an Eulerian path.\hfill$\square$

\begin{notebox}
  The sufficiency proof is constructive: it yields an algorithm. Recursively find a cycle in the graph, remove its edges, and repeat on the remaining graph. This builds Eulerian paths and circuits directly.
\end{notebox}

\subsection*{References}

\begin{itemize}[leftmargin=*, labelsep=0.6em]
  \item M.~Weske. \emph{Business Process Management: Concepts, Languages, Architectures} (3rd ed.). Springer, 2019. \url{http://bpm-book.com}.
  \item Petri Nets World: \url{http://www.informatik.uni-hamburg.de/TGI/PetriNets}
  \item BPMN: \url{http://www.bpmn.org}
  \item Workflow Patterns: \url{http://www.workflowpatterns.com}
  \item WoPeD: \url{http://www.woped.org}
  \item ProM: \url{http://www.promtools.org}
  \item Woflan: \url{http://www.win.tue.nl/woflan}
  \item YAWL: \url{https://yawlfoundation.github.io}
\end{itemize}

